\section{Discussion and conclusions}\label{sec:conclusions}

Two consequences of the theory matter for solver design. First,
certification does not constrain how directions are found: only the final
binary64 row, its bound over the whole block domain and the row that the
solver stores need to be checked, and all of this can be replayed. Second,
the advantage of joint quadratic blocks over their pairs is one of
representation. Pair hulls glued on finitely many moments of a shared
variable accept distributions that no joint distribution reproduces, and they
give the trivial bound zero exactly when the local zero sets alternate often
enough; dense moment relaxations with products that are absent from the model
close the same gaps, and support cuts obtain that strength in the model's own
coordinates. For stars, the most common merged blocks, the support problem
remains exactly solvable in nearly linear time when rows couple the center
with one leaf at a time.

The computations sharpen these conclusions. Certification is a real
safeguard, not a formality: constants taken from samples or from a local
solver were wrong for a large share of the cuts and would have cut off
feasible solutions, including known optima, while every certified cut
replayed. The stored-row check exposed an obstacle that the theory does not
see: presolve aggregates or fixes the variables to which a block refers, so a
separator stated in the original variables loses cuts unless it works in the
presolved space or prevents these reductions, and preventing them weakens
SCIP's own relaxation. On MINLPLib models the cuts were neutral for SCIP, and
SCIP's disabled-by-default separators were stronger, so the cost of the
separator was overhead. On the constructed families, where the theory predicts
a gap, the picture is the opposite. Native SCIP reached the level of glued pair
relaxations once one presolve step was disabled and went no further, and the
certified cuts closed the gap beyond that level and solved every instance. They
did so only with the exact support of each whole row as a first direction,
enough cuts per block, and, when a coupling row binds, a direction search that
finds sloped cuts.

Several questions remain open. Exact support for trees deeper than stars with
edge-local rows would extend the star algorithm to the structures of network
models; Remark~\ref{rem:trees} shows that the rational partition argument
does not extend directly. Proposition~\ref{prop:gain}(ii) reduces the gap of
the glued relaxation in any direction to a concave maximization over
re-splittings, but which directions to merge, and when merging pays off, is
not answered by it. A direction search with the guarantees of
Appendix~\ref{app:separation} and the speed of a solver callback is not
available; on the binding-coupling family the search needed about sixteen cuts
per block to approach the block closure. Finally, an activation rule that
predicts when joint cuts pay off would be needed before any default use; our
structural rule did not achieve this, and on our MINLPLib samples SCIP's own
relaxation already contained what the cuts could add.
